% Shared editorial review version derived from the frozen paper/main.tex.
% Numerical macros, table entries, bibliography and images remain shared with that source.
\section{Introduction}
\label{sec:intro}

Test-Driven Development (TDD)~\cite{beck2002tdd} starts with a failing test, implements
behavior that passes it, and refactors while preserving the tests. Its usefulness depends
on whether those tests represent the intended behavior. An autonomous coding pipeline
can now perform the entire loop: a model translates a specification into tests, generates
code to pass them, and reviews both artifacts before reporting success. The pipeline
produces its own acceptance criteria and its own evidence of meeting them.

Tests and code can nevertheless share a misunderstanding. Practitioners describe tests
that reflect the implementation rather than the requirement as the \emph{ugly mirror},
with tautological tests as an extreme case~\cite{rudolph2008uglymirror,pereira2010ttdd}.
Writing tests first avoids copying an existing implementation but cannot guarantee a
correct interpretation of the task. We measure the \emph{verification gap}: tasks the
pipeline reports as solved whose implementations fail tests withheld from it
(Sec.~\ref{sec:gap}). This is an oracle-relative disagreement; case inspection is needed
to distinguish implementation defects from incomplete requirements or evaluation errors.

We instrument one test-first pipeline and vary which model writes and reviews its tests
and code. Starting with an inexpensive model in every role, we assign stronger models
to test authorship, review, or both while keeping the implementation generator fixed.
Experiments use selected HumanEval and LiveCodeBench problems and three frontier
verifiers. Every configuration retains false accepts. Some estimates favor stronger
verification, and the primary LiveCodeBench reviewer improves externally passing output,
but the false-accept comparisons do not establish a reliable reduction. Residual errors
do not establish that stronger verification cannot help.

RGRBench extends the study to Python packages with stateful and interacting components.
Here stronger review also rejects more work before code generation and reduces correct
output. Acceptance, correctness, and abandonment therefore need to be reported together.
Observed associations with task difficulty and retries do not establish causal effects.

Our contribution is an instrumented comparison of verification roles and pipeline verdicts,
including rejected and abandoned work, across these task settings. We release runners,
prompts, retained results, analyses, and a defect register, while disclosing protocol
deviations and limits of the reconstructed public history.\footnote{Artifact:
\url{https://github.com/integrallis/self-grading-gap}.
RGRBench: \url{https://github.com/integrallis/rgrbench}.}

\section{Background and Related Work}
\label{sec:related}

\paragraph{Prior measurements of self-test error.}
Chen et al. analyze incorrect code accepted by generated tests and correct code rejected
by them~\cite{chen2025revisit}; SFS reports related confusion accounting~\cite{light2025sfs}.
Ahmed et al. compare generated-test-passing repository patches with golden tests and study
refinement~\cite{ahmed2025overfitting}. Our contribution is the executed test-first
instrument, role interventions, and accounting for abandonment, rather than the first
measurement of such disagreement. Zhao et al. separately study tests that pass buggy
implementations and fail their fixes~\cite{zhao2026misguidance}. That code-first mechanism
cannot explain every failure in our initially test-first pipeline.

\paragraph{External checks also need scrutiny.}
EvalPlus strengthens HumanEval evaluation, while LiveCodeBench uses rolling contest tasks
to reduce contamination concerns~\cite{liu2023evalplus,jain2024livecodebench}.
UTBoost and SWE-ABS strengthen repository benchmarks~\cite{yu2025utboost,yu2026sweabs};
BackendForge extends contract-reviewed strengthening to backend services~\cite{guo2026backendforge}.
PatchDiff compares generated and developer patches, but behavioral divergence need not
violate a specification~\cite{wang2025patchdiff}. Conversely, SWE-Bench Pro Verified
repairs overly narrow tests and misleading descriptions~\cite{zheng2026sweproverified}.
Exception-oracle research examines how documentation and code patterns influence expected
behavior~\cite{hossain2026exceptionpatterns}. These studies motivate examining both sides
of a disagreement; their populations and acceptance events differ from ours.

\paragraph{Generated verification can improve.}
CodeT ranks candidates using generated tests~\cite{chen2023codet}; ConVerTest combines
self-consistency, code verification, and execution agreement~\cite{convertest2026}.
SWT-Bench improves accepted-fix precision through generated-test filtering, with a recall
tradeoff~\cite{mundler2024swtbench}. Sol-Ver jointly trains code and test generation and
reduces acceptance of known flawed solutions~\cite{lin2025solver}. CodeHacker improves
rejection using adversarial tests, reference solutions, and human-checked evaluation
infrastructure~\cite{shi2026codehacker}; TCS separately trains test soundness and
adversarial discrimination~\cite{xu2026tcs}. These constructive results concern training,
filtering, or additional reference information. They preclude interpreting our model
substitutions as a general limit on verification; their error rates also use different
denominators.

\paragraph{Input and training regimes matter.}
Mathews and Nagappan supply tests during generation~\cite{mathews2024tdd}; TiCoder uses
interactive feedback to clarify intent~\cite{fakhoury2024ticoder}. TENET uses
developer-written tests, while TDFlow studies supplied and generated suites in repository
workflows~\cite{sun2025tenet,tdflow2025}. MetaGPT studies multi-agent software
development~\cite{hong2024metagpt}. ReVeal trains explicit self-verification and reports
useful iterative refinement~\cite{jin2025reveal}. These settings differ from our prompted
reviewers and caution against interpreting observational retry groups as a general
refinement penalty.

\paragraph{Verdicts, selection, and feedback.}
CodeJudgeBench compares model judgments with execution~\cite{jiang2025codejudgebench};
Jin and Chen study false rejection of correct code~\cite{verifyfailures2025}.
SWE-RM evaluates discrimination, calibration, and selection on fixed
trajectories~\cite{shum2025swerm}. RETRACE reconstructs an issue from a patch without
the original issue, then reconciles and revises~\cite{li2026retrace}. Improved final
resolution alone does not measure verdict reliability. Our judge intervention likewise
changes generation. Imperfect-verifier resampling~\cite{stroebl2026resampling} and
General AgentBench's difference between oracle pass@$K$ and model selection~\cite{li2026generalagentbench}
address related but distinct endpoints.

\paragraph{Specifications and complementary evidence.}
The oracle problem is established~\cite{barr2015oracle}. VERINA separately evaluates
code, specifications, and proofs~\cite{ye2026verina}; proving a formal contract leaves
its fidelity to intended requirements to be assessed. Sam-Bodden's context
ablation~\cite{sambodden2026context} varies code representation after fixing localization
and checks external SWE-bench outcomes. It is complementary methodology, not a replication
of the present tests, interventions, or disagreement rates.


\section{The Verification Gap}
\label{sec:gap}

\subsection{Definitions}
For task $i$, the pipeline authors tests $T_i$ from a natural-language specification,
generates implementation $y_i$, and records a self-verdict $s_i\in\{0,1\}$.
An independent oracle verdict $o_i\in\{0,1\}$ comes from tests withheld from the
pipeline. Acceptance requires successful final execution of the self-tests;
rejection can also result from a judge gate or exhausted retries, including
abandonment before any implementation exists.

The verification gap is $G=\{i:s_i=1\wedge o_i=0\}$, the oracle-relative
\emph{false accepts}. We report its count and its conditional rate,
$|G|/|\{i:s_i=1\}|$, among claimed successes. False rejects are
$R=\{i:s_i=0\wedge o_i=1\}$. True accepts and true rejects have
$s_i=o_i=1$ and $s_i=o_i=0$, respectively. Reporting all four cells distinguishes
an accurate decision from a policy that rarely produces or accepts code.

Our oracles differ in distance from the prompts. HumanEval's original tests
overlap visible worked examples~\cite{chen2021humaneval}; HumanEval\texttt{+}
substantially expands that suite~\cite{liu2023evalplus}. LiveCodeBench provides
private tests withheld from its problem statements~\cite{jain2024livecodebench}.
Every false accept is a disagreement with the selected oracle. An omitted
self-test, incorrect expected value, implementation defect, underspecified
requirement, or oracle/adapter error can cause that disagreement. The definition
alone does not establish a specification violation or shared misunderstanding.
The RGRBench example in Sec.~\ref{sec:results-scale} makes this distinction concrete.

\subsection{Pipeline and review inputs}
\label{sec:prompts}
The RED test author sees the task's specification and worked examples. Its prompt
requires exact signatures, specified behavior, recalculated expected values,
ordinary pytest, and a canonical test for each worked example. The RED judge sees
the specification and proposed tests, checks example coverage, assertion fidelity,
invented constraints, and coverage, and returns a score plus
\texttt{ADD:}/\texttt{REMOVE:}/\texttt{FIX:} directives. A deducted error can still
leave the score above the acceptance threshold. Failed attempts receive feedback
and the preceding attempt.

GREEN implements the generated tests, with the task title, available implementation
hints, retained implementation files, and retry feedback. For HumanEval and LCB,
these hints include the full original specification; RGRBench supplies generic
implementation instructions instead. Its prompt prohibits
hardcoding and returns implementation code as an escaped JSON field. The GREEN
judge receives the generated tests, implementation, and task title rather than the
full original specification. It predicts per-test outcomes and flags test-gaming.
If its score clears the gate, execution checks the implementation against the
generated suite and overrides an optimistic prediction. VERIFY runs the full suite;
failures can trigger further GREEN attempts. A failure analyst returns textual fix
directives, never code. Thus stronger judges operate under a particular
test-centered contract, not an unrestricted specification-aware review.
% trace: src/vgap/pipeline/templates/agents/code/python_generate_from_tests.j2;
% src/vgap/pipeline/templates/judge/code_against_tests.j2; coordinator.py

\begin{figure*}[t]
\centering
\includegraphics[width=\linewidth]{figures/fig1_loop.png}
\caption{The pipeline authors tests $T$ and code $y$, grades
itself by running $T(y)$, and emits self-verdict $s_i$; the oracle sees $y$ but never $T$.
The verification gap is the set of problems where the two verdicts disagree in the pipeline's
favor.}
\label{fig:loop}
\end{figure*}


The artifact retains the instrument, prompt templates, render-contract checks,
canonical-example anchoring variants, and the full defect register. These details
support inspection; the core measurement does not depend on a supplement being
available to reviewers.

\subsection{Independent measurement and repeated runs}
\label{sec:protocol}
The harness records the executed pipeline verdict and captures generated artifacts.
It then runs the benchmark oracle on the implementation outside the feedback loop:
EvalPlus, the LiveCodeBench evaluator, or RGRBench suites through the constrained
adapter described below. Oracle tests and results never enter generation or review.
The records support confusion matrices, acceptance, oracle success, and retained-call
cost accounting.

Oracle strength changes the measurement: the primary all-weak HumanEval control
has $3.0$ false accepts per 30-problem run against the original oracle and $4.0$
against HumanEval\texttt{+}. Byte-identical configurations also vary across runs:
the control disagrees with itself on $6/30$ HumanEval tasks and $9/30$
LiveCodeBench tasks across five runs, with maximum pairwise pass-total differences
of $4$ and $3$. The protocols combine statistical tests with minimum per-run
effect sizes of $4$ for oracle success, $2$ for false accepts, $2$ for false
rejects, and $3$ for total verdict error. These thresholds are neither confidence
bounds nor equivalence margins; some control spreads exceed them. Section~\ref{sec:design}
distinguishes the original decision rules from the post hoc reporting audit.
% trace: experiments/exp006_model_asymmetry/results/ANALYSIS.md;
% experiments/exp007_lcb_asymmetry/results/ANALYSIS.md


\section{Study Design}
\label{sec:design}

The study asks four questions.

\noindent\textbf{RQ1 (phenomenon).} How large is the verification gap, against visible-adjacent,
strong, and hidden oracles, and how does it break down into false accepts and false rejects?

\noindent\textbf{RQ2 (asymmetry).} Does a frontier model in the verification roles (as test author,
as LLM-as-a-judge reviewer, or both) close the gap?

\noindent\textbf{RQ3 (mechanism).} When a stronger verifier does help, does the improvement come
from its \emph{verdicts} (better gatekeeping) or from its \emph{feedback} (better generation)?

\noindent\textbf{RQ4 (scale).} Do the size of the gap and its resistance to verifier strength persist
when moving from single-function algorithms to larger, application-scale programs that involve
multiple concerns, and how does this gap change with increasing task complexity or as the loop
itself is refined?

\paragraph{Experiments.} The model experiments use a single asymmetry design with four conditions
(\emph{cells}): a weak model W (\texttt{gpt-4o-mini}) generates code in every cell; a frontier
verifier S replaces W in the test-author role, the judge role (covering all three judge duties), or
both; the control condition uses W in every role. We assign the S role to three verifiers from
three different providers---GPT-5.6 Terra, Gemini-2.5-Pro, and Claude-4.5-Sonnet\footnote{\texttt{gpt-4o-mini} is
OpenAI's GPT-4o mini. The verifiers are OpenAI's \texttt{gpt-5.6-terra}, Google's \texttt{gemini-2.5-pro}
(via Vertex AI), and Anthropic's \texttt{claude-sonnet-4-5}. All models are called through a LiteLLM
shim.}---and run this design on both benchmarks for each verifier. $N{=}5$ runs per cell, $30$
problems per run, model identity is fixed within each assigned role, and judge threshold and retries are held
constant. The single-function runners can switch late generation attempts to a second client
for the same model at a different temperature. The protocols record advance hypotheses,
but the public release history was reconstructed: original chronology depends on private
commits (\texttt{PROVENANCE.md}). Exp011 documents collection and retention changes during
the runs; these limit a strictly confirmatory interpretation. The primary pair uses GPT-5.6 Terra: exp006 runs a HumanEval subset that
includes previously unstable problems, scored by both EvalPlus oracles, and exp007 repeats this
on contamination-controlled data---all $19$ easy plus the first $11$ medium LeetCode problems from
LiveCodeBench release~v6 with contest dates $\geq$~2025-01-01, chosen by a deterministic rule fixed
in advance after a pilot identified the weak model's measurable range, and scored by the private LCB
suites. The exp007 protocol records predictions of fewer false rejects and no detectable
false-accept contrast, following exp006. The exp011 replication then reruns
both benchmark designs with the two non-OpenAI verifiers, Gemini-2.5-Pro and Claude-4.5-Sonnet, to
test whether the soundness null depends on which frontier model is used (Sec.~\ref{sec:results-asym}).
% trace: experiments/exp006_model_asymmetry/README.md, exp007_lcb_asymmetry/README.md,
% exp011_verifier_generalization/PRE_REGISTRATION.md (pre-registrations);
% experiments/pilots/README.md (band-location pilot)

\paragraph{The application-scale experiment.} A third experiment (exp010) moves to
RGRBench's $89$ packages. The pipeline receives only each package's requirements document
(user-story acceptance criteria derived solely from the tests, with the tests themselves never
shown) and is graded by the hidden, hand-verified suite through a non-cheating adapter
(Sec.~\ref{sec:results-scale}). The design keeps the judge-role contrast: a \emph{control} cell (weak model in every role) versus a \emph{strong-judge}
cell (GPT-5.6 Terra in all judge duties), with $N{=}3$ runs per package per cell ($534$ runs). Since
RGRBench labels each package by difficulty tier and the harness records whether a claim was made on
the first try or only after retries, these runs describe associations with complexity (H-D4) and additional attempts (H-D2).
Neither difficulty tier nor retry count is randomized.
% trace: experiments/exp010_flagship/README.md (pre-registration); results/plan.json

\paragraph{Statistical protocol and reporting correction.} For oracle-pass and false-accept
contrasts in the single-function experiments, the original analyzers take each problem's
majority outcome over five runs and apply exact McNemar tests, alongside practical thresholds
on mean per-run differences. We call these combined requirements the \emph{decision bar}.
Majority outcomes and mean changes are different estimands. Failing this bar does not establish
absence of an effect. The original scripts omit paired tests for false rejects and total
verdict error, and do not apply the protocols' larger-control-spread adjustment; these are
analysis deviations. Reported $p$-values are nominal, without multiplicity correction.

The original HumanEval summaries combine base-oracle pass and false-reject counts with
HumanEval\texttt{+} false accepts. Table~\ref{tab:cells} corrects this to use
HumanEval\texttt{+} throughout, recomputed from retained scores. The original summaries
remain available. A separately labeled post hoc audit provides complete matrices for both
oracles, the missing paired tests, and problem/package-cluster bootstrap intervals that
retain repeated runs together (\texttt{experiments/review\_audit/}). These supplemental
analyses do not retroactively become pre-registered confirmations.

\paragraph{Telemetry.} Retained successful LLM calls are logged with token counts, cost, and a problem, cell,
or run tag. This telemetry enables the mechanism analysis described in
Sec.~\ref{sec:results-mech}, and it is included with the artifact. The primary configurations use code-generation temperature $0.7$,
weak/frontier test-author temperatures $0.7/1.0$, judge temperature $1.0$,
threshold $\tau{=}0.70$, and $3$ RED/GREEN attempts. Late single-function generation
can use the same model at temperature $1.0$. Self-tests have a $30$\,s timeout;
oracle limits are harness-specific (LCB parameter $6$\,s; RGRBench $120$\,s).
Exp013 varies the threshold. Provider handling and adapter token limits differ.
% trace: per-run usage.jsonl files; totals in each asym_summary.json



\paragraph{Threats addressed by design.} HumanEval is the field's best-known benchmark and
predates the training data of all current models, so we consider it within-distribution. The
LiveCodeBench window serves as the off-distribution check, since its problems were released more
than a year after the weak model's training data. In the audited primary runs, we checked all $7{,}962$ logged prompts from both experiments and confirmed that rendered
benchmark identity is conveyed through the statement and worked examples, without a
benchmark-name or problem-number label; generated tests, code, and feedback also appear in prompts. A date after the generator's reported cutoff reduces contamination concerns but does not
prove absence of exposure, particularly for later verifiers (Sec.~\ref{sec:threats}). Both problem sets are selected for measurability (with increased hardness; within the
weak-model band), so per-run counts do not represent benchmark rates and are never reported as
such.

\section{Results}
\label{sec:results}

\subsection{RQ1: the gap on the corrected instrument}
\label{sec:results-gap}

Table~\ref{tab:cells} presents the full result matrix and error breakdown. On the all-weak control, the gap appears on both benchmarks and both sides.
For each 30-problem run, HumanEval\texttt{+} has $\AuditHEControlPass$ oracle passes,
$\AuditHEControlFA$ false accepts, and $\AuditHEControlFR$ false rejects in control.
LiveCodeBench has $15.2$ oracle passes, $4.8$ false accepts ($29\%$ of claims), and $3.4$
false rejects. The LCB problem dates are after the generator's reported cutoff, and private
tests are not included in prompts. These controls reduce exposure concerns without proving
that every problem is novel to every model. Both error labels are relative to the chosen
oracle and the pipeline's complete set of acceptance gates.
% trace: exp006 ANALYSIS.md control_rerun; exp007 ANALYSIS.md control

\begin{table*}[t]
\centering
\caption{Complete result matrix: per-run means over $N{=}5$ runs of 30 problems (oracle-consistent HumanEval\texttt{+}
values from the post hoc reporting correction). FA = false accepts (HumanEval: strong oracle; LCB: hidden
oracle); FR = false rejects. The implemented decision bar uses a fixed minimum effect size (pass 4, FA 2, FR 2)
\emph{and} exact McNemar $p{<}.05$; $^{\dagger}$ = past that fixed effect-size threshold,
$^{\ast}$ = clears the full bar.}
\label{tab:cells}
\small
\begin{tabular}{llrrrr}
\toprule
Benchmark & Cell & Oracle pass & FA & FR & tests/run \\
\midrule
HumanEval subset-30 & control (W everywhere) & \AuditHEControlPass & \AuditHEControlFA & \AuditHEControlFR & 84 \\
(strong oracle)     & S authors tests & \AuditHETestPass & \AuditHETestFA & \AuditHETestFR & 89 \\
                    & S judges & \AuditHEJudgePass & \AuditHEJudgeFA & \AuditHEJudgeFR$^{\dagger}$ & 105 \\
                    & S authors + judges & \AuditHEBothPass & \AuditHEBothFA & \AuditHEBothFR$^{\dagger}$ & 101 \\
\midrule
LiveCodeBench lcb30 & control (W everywhere) & 15.2 & 4.8 & 3.4 & 76 \\
(hidden oracle,     & S authors tests & 16.4 & 2.8 & 2.2 & 140 \\
post-cutoff)        & S judges & 20.2$^{\ast}$ & 2.4$^{\dagger}$ & 0.0$^{\dagger}$ & 182 \\
                    & S authors + judges & 21.4$^{\ast}$ & 3.0 & 0.0$^{\dagger}$ & 194 \\
\bottomrule
\end{tabular}
\end{table*}



\subsection{RQ2: frontier models in the verification roles}
\label{sec:results-asym}

\paragraph{No primary false-accept contrast meets the decision bar.} On HumanEval, mean
false-accept counts change from $4.0$ in control to $3.2/4.4/3.4$ under a strong test author,
judge, or both (all nominal $p{=}1.0$). On LiveCodeBench, they decrease from $4.8$ to
$2.8/2.4/3.0$, with nominal $p \in [0.125, 0.375]$. These estimates include potentially
useful reductions, especially for the strong judge. The primary tests do not support the
specified positive claim; equally, their null results do not establish invariance. False
accepts persist in every measured condition.
% trace: exp006/exp007 results/ANALYSIS.md (FA contrasts)

\paragraph{Residual errors under two more frontier verifiers.} Exp011 repeats the design
with Gemini-2.5-Pro and Claude-4.5-Sonnet. CodeJudgeBench reports strong results for
Gemini-2.5-Pro and the earlier Claude-4-Sonnet, but did not evaluate the Sonnet 4.5 model used
here~\cite{jiang2025codejudgebench}. No false-accept contrast meets the specified bar
(Table~\ref{tab:verifiers}); several point estimates decline. False-reject changes are not
uniform across providers: a strong judge does not eliminate them in every replication.
Collection-time validity and rerun changes are documented in exp011's deviation log and
must accompany interpretation of these retained runs.
% trace: experiments/exp011_verifier_generalization/results/*/ANALYSIS.md; DEVIATIONS.md

\begin{table*}[t]
\centering
\small
\caption{False accepts per 30-problem run (mean, $N{=}5$) under three frontier verifiers, each
swapped into the test-author (\emph{S tests}), judge (\emph{S judge}), or both roles against the
fixed weak generator. For every verifier, on both benchmarks, no strong-cell contrast clears the
decision bar. This is not an equivalence test against the all-weak control.
GPT-5.6 Terra is exp006/exp007; Gemini-2.5-Pro and Claude-4.5-Sonnet are the exp011
replication, with collection deviations disclosed (Sec.~\ref{sec:results-asym}). HumanEval is scored with HumanEval\texttt{+}.}
\label{tab:verifiers}
\begin{tabular}{lrrrr}
\toprule
Frontier verifier S & control & S tests & S judge & S both \\
\midrule
\multicolumn{5}{l}{\emph{HumanEval}\texttt{+}} \\
\quad GPT-5.6 Terra     & $4.0$ & $3.2$ & $4.4$ & $3.4$ \\
\quad Gemini-2.5-Pro    & $4.8$ & $3.2$ & $4.0$ & $4.8$ \\
\quad Claude-4.5-Sonnet & $3.8$ & $2.8$ & $3.4$ & $4.2$ \\
\midrule
\multicolumn{5}{l}{\emph{LiveCodeBench} (hidden, post-cutoff oracle)} \\
\quad GPT-5.6 Terra     & $4.8$ & $2.8$ & $2.4$ & $3.0$ \\
\quad Gemini-2.5-Pro    & $3.8$ & $1.8$ & $3.4$ & $2.6$ \\
\quad Claude-4.5-Sonnet & $3.6$ & $2.0$ & $3.4$ & $3.6$ \\
\bottomrule
\end{tabular}
\end{table*}



\paragraph{The primary frontier judge reduces observed false rejects.} On LiveCodeBench,
both primary frontier-judge cells have zero false rejects in all ten runs (control: $3.4$).
On HumanEval\texttt{+}, false rejects fall from $\AuditHEControlFR$ to $\AuditHEJudgeFR$
(judge) and $\AuditHEBothFR$ (both). These are descriptive reductions; the completed
majority-outcome tests do not establish a false-reject effect under the full decision rule.
% trace: experiments/review_audit/results/ANALYSIS.md

\paragraph{The primary frontier judge increases LiveCodeBench oracle passes.} Oracle
passes rise from $15.2$ to $20.2$ (judge; $\Delta{+}5.0$, nominal $p{=}.0312$) and $21.4$
(both; $\Delta{+}6.2$, nominal $p{=}.0156$), meeting the original decision bar. Neither
corresponding HumanEval\texttt{+} contrast meets that bar. These nominal tests are not
adjusted for the multiple comparisons. The post hoc audit also computes tests omitted by
the original analyzer, so these should not be described as the only detectable outcomes.
% trace: exp007 results/ANALYSIS.md; experiments/review_audit/results/ANALYSIS.md

\paragraph{No pass-rate benefit is established for a frontier test author.} The
strong-test-author cell has higher mean oracle passes than control on both benchmarks, but
neither primary comparison meets the decision bar. FA and FR point estimates also decrease.
The result is a failure to establish the specified benefit, not proof that test authorship
quality has no effect.

\subsection{RQ3: telemetry is consistent with a feedback benefit}
\label{sec:results-mech}

Canonical tests average $75$ per run in the LCB cells. Total test volume changes from about
$76$ in control to $182$--$194$ in the frontier-judge cells, while cycle completions rise from
about $17/30$ to $23$--$24/30$. A frontier test author also produces larger suites ($140$
tests per run), without a pass-rate contrast meeting the decision bar.
% trace: exp007 results/ANALYSIS.md; retained launch logs

These observations are consistent with stronger feedback helping generation. They do not
identify feedback as the cause: the judge intervention jointly changes scores, directives,
retry paths, test volume, and the selected outputs. The strong-test-author condition also
retains iterative review, so it is not a one-shot comparator. A feedback-only versus
verdict-only ablation on matched candidates would be needed to separate these mechanisms.

\subsection{RQ4: oracle disagreement at application scale}
\label{sec:results-scale}

In the RGRBench control, $65$ of $261$ valid runs end in a claimed success, of which
$41.5\%$ fail the oracle. The strong-judge cell claims $50$ successes among $251$ valid runs,
with a conditional gap of $42.0\%$ (Table~\ref{tab:scale}). The planned $267$ runs per cell
reduce to these valid totals because evaluations without usable oracle verdicts are excluded:
$21$ exclusions follow adapter import failures at test collection, and one control run has a
harness environment error. A post hoc missing-verdict sensitivity analysis is included with
the artifact. It does not resolve possible oracle or adapter errors among valid cases.
% trace: exp010 results/ANALYSIS.md; experiments/review_audit/results/ANALYSIS.md

The table preserves nominal run-level Wilson intervals. These treat repeated package-runs
as independent and understate uncertainty when their errors cluster. A paired bootstrap that
resamples whole packages gives a strong-minus-control conditional-gap difference of
$\AuditRGRSharedDelta$ percentage points, with interval $\AuditRGRSharedCI$. This allows
substantial improvement and harm; the data do not establish equivalence. Cross-benchmark
point estimates are descriptive because tasks, oracle contracts, and adapters differ.
% trace: experiments/review_audit/results/paper_numbers.tex

RGRBench grades candidates through a generated adapter constrained to expose their APIs
under the oracle's import names. The selected calibration fixtures have no observed adapter
failures, but this is not a zero error bound (Sec.~\ref{sec:threats}). The example below
also demonstrates that oracle failure need not mean violation of the supplied requirement.

\begin{figure*}[t]
\centering
\includegraphics[width=\linewidth]{figures/fig6_requirements_oracle.png}
\caption{A requirements--oracle mismatch (RGRBench's \texttt{numbers\_to\_words}).
AC-5.1 specifies the error message but leaves the exception type unspecified. The self-test
accepts any \texttt{Exception}, which the code raises; the oracle additionally requires
\texttt{ValueError}. The resulting false accept is relative to the oracle, not demonstrated
violation of the visible text. Code snippets are simplified; full excerpts appear in
the released long version.}
\label{fig:misalignment}
\end{figure*}

\begin{table*}[t]
\centering
\caption{RGRBench (application scale): the self-verification confusion matrix and the
false-accept analyses, over $N{=}3$ runs of $89$ packages ($534$ runs). Control = weak model
everywhere; strong-judge = frontier model in all judge roles. TA/FA/FR/TR are counts over
valid runs (accept/reject $\times$ oracle pass/fail); oracle-passing solutions $=$ TA${+}$FR.
Conditional FA is FA as a fraction of claimed successes (TA${+}$FA); tier and attempt rows
restrict that quantity to each subset. Brackets are nominal $95\%$ run-level Wilson
intervals; package-cluster uncertainty is reported in the text.}
\label{tab:scale}
\small
\begin{tabular}{lcc}
\toprule
 & control & strong-judge \\
\midrule
valid runs                           & 261 & 251 \\
true accepts (TA)                    & 38  & 29  \\
false accepts (FA)                   & 27  & 21  \\
false rejects (FR, H-D3)             & 14  & 0   \\
true rejects (TR)                    & 182 & 201 \\
oracle-passing solutions (TA${+}$FR) & 52  & 29  \\
\midrule
claimed successes (TA${+}$FA)        & 65 & 50 \\
conditional FA (H-D1) & 0.415 [0.304, 0.537] & 0.420 [0.294, 0.558] \\
\quad beginner tier (H-D4) & 0.31 & 0.38 \\
\quad intermediate tier & 0.54 & 0.50 \\
\quad advanced tier & 0.86 & 0.75 \\
first- / multi-attempt FA (H-D2) & 0.35 / 0.64 & 0.34 / 0.67 \\
\bottomrule
\end{tabular}
\end{table*}

\paragraph{Higher observed gap rates on harder packages.} Conditional false accepts rise
across difficulty tiers: $31\% \to 54\% \to 86\%$ in control,
and $38\% \to 50\% \to 75\%$ under the strong judge. The hardest control tier has only seven
claimed successes, six rejected by the oracle. These small, selected denominators and
nonrandom difficulty assignments limit any causal or general scaling interpretation.
% trace: exp010 results/ANALYSIS.md (H-D4 rows)



\paragraph{Additional attempts are associated with higher gap rates.} Claims reached after
multiple GREEN attempts have higher conditional false accepts than first-attempt claims:
$64\%$ versus $35\%$ in control, and $67\%$ versus $34\%$ with the strong judge. Retry count
is a consequence of earlier failures; difficult cases are more likely to retry. The design
therefore cannot distinguish selection from a harmful effect of refinement, and it does not
provide a causal replication of the refinement intervention in Ahmed et al.~\cite{ahmed2025overfitting}.
% trace: exp010 results/ANALYSIS.md (H-D2 rows)

\paragraph{Acceptance and correct output must be considered together.} Although the
conditional gap estimates are close ($41.5\%$ versus $42.0\%$), the strong judge reduces
oracle-passing outputs from $52$ to $29$ and true accepts from $38$ to $29$. It rejects the
self-authored suite before code generation on $189$ of $251$ valid runs; only $62$ reach
GREEN, compared with $258$ under the weak judge. Zero false rejects consequently does not
imply improved recovery of correct code. The intervention changes which implementations
exist and which are accepted, so this is a comparison of pipeline policies, not judge
accuracy on fixed candidates.
% trace: exp010 results/ANALYSIS.md; result.json green_attempts

\paragraph{Threshold follow-up: different operating points, incomplete matching.} Exp013
varies the threshold shared by the RED and GREEN judges. The original new points are
$\tau\in\{0.40,0.55,0.85\}$, with $N{=}1$ per package; $0.70$ reuses exp010. A later
$\tau{=}0.50$ follow-up uses $N{=}3$ (Table~\ref{tab:tausweep}). That point is outside the
original matching candidate set, and no separately frozen selection amendment is present;
we treat it as adaptive follow-up rather than confirmatory matching.
% trace: exp013 PRE_REGISTRATION.md; results/ANALYSIS.md

Acceptance at $0.50$ is $31.5\%$, versus $24.9\%$ in control, so throughput was not equalized.
Its conditional gap is $47.5\%$, versus $41.5\%$ in control. A post hoc paired package
bootstrap gives a difference of $\AuditRGRFollowupDelta$ percentage points with interval
$\AuditRGRFollowupCI$, supporting neither superiority nor equivalence. At $0.40$, both
acceptance and the observed gap are higher; at $0.85$, acceptance is low and the gap estimate
has wide uncertainty. The rates are not monotonic at every intermediate point. The sweep
shows sensitivity to the operating threshold, but a weak-judge threshold curve and matched
candidate evaluation would be needed to compare discrimination at similar throughput.
% trace: experiments/review_audit/results/paper_numbers.tex; exp013 results/ANALYSIS.md

\begin{table*}[t]
\centering
\caption{RED/GREEN threshold sweep on the strong-judge RGRBench cell (exp013).
The adaptive $\tau{=}0.50$ follow-up has $N{=}3$; other new points have $N{=}1$;
$\tau{=}0.70$ reuses exp010 ($N{=}3$). Acceptance is not matched to control.
Brackets are nominal $95\%$ run-level Wilson intervals.}
\label{tab:tausweep}
\small
\begin{tabular}{lccc}
\toprule
cell & accept\% & FA\% [$95\%$] & RED-gate\% \\
\midrule
control (weak), $\tau{=}0.70$      & $24.9$ & $41.5$ [$30.4, 53.7$] & $1.1$  \\
strong, $\tau{=}0.85$              & $7.4$  & $33.3$ [$9.7, 70.0$]  & $91.4$ \\
strong, $\tau{=}0.70$              & $19.9$ & $42.0$ [$29.4, 55.8$] & $75.3$ \\
strong, $\tau{=}0.55$              & $20.0$ & $52.9$ [$31.0, 73.8$] & $72.9$ \\
strong, $\tau{=}0.50$ ($N{=}3$)    & $31.5$ & $47.5$ [$36.9, 58.3$] & $62.6$ \\
strong, $\tau{=}0.40$              & $43.0$ & $67.6$ [$51.5, 80.4$] & $47.7$ \\
\bottomrule
\end{tabular}
\end{table*}



\section{Discussion}
\label{sec:discussion}

\paragraph{Residual error and uncertainty.}
Every configuration retains false accepts, but failure of the original joint
decision rule does not establish invariance or equivalence. Several point
estimates favor stronger verification. False-accept counts and error among
accepted outputs also answer different questions because interventions change
both acceptance and the candidates produced.

\paragraph{Feedback, selection, and the review contract.}
The primary LiveCodeBench gains are consistent with useful reviewer feedback,
but role reassignment simultaneously changes test acceptance, code review, failure
analysis, and subsequent candidates. It does not isolate coaching from selection.
The RGRBench reduction in correct output and the Gemini judge-only false-reject
results limit generalization. Deployment recommendations additionally require
configuration-specific cost and outcome comparisons.

The GREEN judge evaluates code against generated tests without the full original
specification, while the RED rubric can accept a suite with a deducted error.
An internal reviewer given the specification could still identify omissions;
the results do not demonstrate that this is impossible. This restriction connects
the study to the oracle problem and imperfect-verifier inference
scaling~\cite{barr2015oracle,stroebl2026resampling}.

\paragraph{What remains unmeasured.}
A direct next baseline would freeze candidates and compare specification-aware
review predictions against the same external oracles. A separate randomized
ablation would withhold feedback while preserving verdicts. Neither has been
run. Fixed-trajectory review and reconstruction-based checking provide relevant
precedents~\cite{shum2025swerm,li2026retrace}.

The requirements--oracle mismatch also prevents treating every held-out failure
as a semantic defect. Verification and validation are distinct~\cite{boehm1984vv};
intent clarification~\cite{fakhoury2024ticoder} and blinded adjudication against
both requirements and oracle contracts are plausible follow-ups, not measured
remedies here. For builders, self-test success is evidence whose reliability
must be evaluated on the intended tasks and acceptance policy. Independent tests,
checked specifications, and differential or metamorphic checks can add evidence,
each with its own limitations.


\section{Threats to Validity}
\label{sec:threats}

\paragraph{Construct and oracle adequacy.} All correctness labels are relative to the stated
benchmark oracle. Incomplete tests can miss defects, while an incorrect or underspecified oracle
contract can reject behavior allowed by the prompt. HumanEval's base and plus suites are distinct
oracles; the post hoc reporting correction uses complete matrices from each rather than combining
base passes and false rejects with plus false accepts. On RGRBench, the mutation kill rate
characterizes detection of selected mutations, not the probability that the oracle's expected
behavior matches the requirements. Classifying surviving mutants does not provide a semantic
error bound. Coverage and mutation evidence also depend on the testing regime, as recent
replication work distinguishes tests for correct code from tests intended to expose existing
bugs~\cite{zhao2026coverage}. Positive self-verdicts require passing execution; negative verdicts additionally
include judge rejection, exhausted retries, or failure to produce an implementation. Thus false
rejects are not exclusively cases of faulty tests vetoing otherwise correct code.
% trace: experiments/review_audit/results/ANALYSIS.md (oracle-consistent matrices);
% src/vgap/pipeline/subgraph.py; experiments/exp010_flagship/run_flagship.py

\paragraph{Adapter calibration.} The calibration exercises model-written adapters on
behavior-preserving transformations of reference implementations. It observes no adapter errors
in $13$ retained fixtures, after excluding $5$ fixtures whose deterministic adaptation failed.
These selected cases preserve much of the reference structure and do not cover arbitrary
candidate APIs. The recorded calibration error is a fraction of tests failed, whereas the study's
binary endpoint fails a candidate if any oracle test fails. Zero observed calibration errors
neither gives a zero population error bound nor excludes adapter-caused disagreements in the
experimental runs. Adapter and specification mismatches remain possible contributors to the
application-scale measurements.
% trace: experiments/exp010_flagship/calibration/calibration.json;
% experiments/exp010_flagship/calibrate_adapter.py

\paragraph{Selection and repeated packages.} Model and threshold changes alter both acceptance
and the candidates reaching a usable oracle verdict. Restricting a comparison to packages valid
in both cells does not match accepted candidates; similar throughput would not do so either.
The threshold experiment changes RED and GREEN together, and the $\tau{=}0.50$ follow-up is
adaptive and does not exactly match control acceptance. These are comparisons of pipeline
policies, not estimates of reviewer accuracy on a fixed candidate population. Nominal Wilson
intervals treat package-runs as independent; repeated runs on the same package can be correlated.
The supplementary package bootstrap retains repeats and pairing and leaves broad uncertainty
about the conditional-error contrasts. It describes resampling of these selected packages,
not representative sampling of software projects. Three new threshold points have only one
run per package, limiting interpretation of the operating curve. Retry and difficulty-tier
comparisons are also observational: retrying tasks can be harder before the retry occurs.
% trace: experiments/exp013_tau_sweep/PRE_REGISTRATION.md; run_tau_sweep.py;
% experiments/review_audit/results/ANALYSIS.md (package bootstrap)

\paragraph{Registration and collection deviations.} The authors retained protocols intended to
precede collection, but the public history was reconstructed for release; historical run SHAs
refer to a private working history. The released Git graph alone does not independently establish
the original registration chronology. In exp011, partial runs were deleted and repeated after
provider outages, using a nonempty-implementation inclusion threshold that was subsequently
relaxed. This differs from the protocol's treatment of terminal API errors as failed runs and can
also select on legitimate pipeline abandonment. Without the discarded attempts, the resulting
selection and cost effects cannot be fully reconstructed. Cross-provider findings therefore
refer to retained runs, with this additional limitation.
% trace: PROVENANCE.md;
% experiments/exp011_verifier_generalization/PRE_REGISTRATION.md; DEVIATIONS.md

\paragraph{Conclusion validity.} The original tests dichotomize each problem's repeated outcomes
by majority, while effect floors use mean per-run counts; these are different estimands. Some
registered false-reject and total-error tests and the control-spread adjustment were omitted by
the original runners. The review supplies completed tests and clustered intervals as explicitly
post hoc analyses, with nominal, unadjusted $p$-values. They must not be relabeled as prospective
confirmations. Neither failure of a significance test nor failure to exceed a practical floor
establishes an effect of zero. With small selected task sets, repeated outcomes, and numerous
contrasts, meaningful benefits and harms remain compatible with several comparisons.
% trace: experiments/review_audit/results/ANALYSIS.md;
% docs/STATISTICAL_REVIEW.md

\paragraph{Scope and contamination.} The main interventions keep one implementation-generator
model fixed in one pipeline. All three verifiers are evaluated on the two single-function
subsets; only GPT-5.6 Terra is evaluated on RGRBench. Model aliases and available historical
metadata do not fully identify immutable served snapshots. The LiveCodeBench window postdates
the weak model's documented cutoff and reduces a known exposure risk, but does not establish
that task ideas or variants were unseen. Later verifiers may have encountered the tasks;
exposure could change tests, acceptance, feedback, and candidate selection, so its effect on
conditional error has no established direction. Differences between benchmarks also combine
changes in problem selection, specifications, and scoring contracts with application scale.
The results should not be treated as benchmark-wide rates or a general complexity law.

\paragraph{Artifact and execution limits.} The release supports reanalysis of stored outcomes,
with incomplete historical provider snapshots, request settings, and some execution traces;
exact API replay is not established. Logged costs estimate retained successful calls and can
omit discarded or failed attempts; offline adapter grading also contributes to study spend.
Generated code runs in temporary directories under subprocess timeouts, not a hardened sandbox
with network, memory, and filesystem isolation. Timeout policies differ between self-testing
and the benchmark harnesses. Fresh replication should record these settings and use appropriate
execution isolation.

\paragraph{Anchoring.} The artifact includes canonical-example anchoring variants, but these do
not provide a false-accept comparison sufficient to establish that anchoring cannot help. The
case studies show that covering worked examples can coexist with a held-out failure. A fully
instrumented anchoring ablation with repeated runs remains a separate experiment.

\section{Conclusion}
\label{sec:conclusion}

A self-testing pipeline can accept implementations that fail held-out tests. We measure
that disagreement across selected single-function problems and RGRBench packages, and show
why acceptance, oracle success, and abandonment must be reported together. Frontier
verification leaves residual false accepts in every measured condition, while some
interventions improve oracle success and several false-accept estimates decrease. The
experiments do not establish that verifier strength cannot help or that error rates are
invariant. At application scale, stricter gating reduces correct output, and requirements,
oracles, and adapters add measurement uncertainty.

For builders, passing self-authored tests is useful evidence with an empirically measurable
error rate, not a correctness guarantee. Independent evaluation and explicit verdict
accounting remain necessary to assess that evidence. The released instrument and post hoc
reporting corrections provide a basis for stronger tests of specification-aware review,
feedback mechanisms, and oracle fidelity across more pipelines and generators.
